
\documentclass{article}
\usepackage{colortbl}
\usepackage{makecell}
\usepackage{multirow}
\usepackage{supertabular}

\begin{document}

\newcounter{utterance}

\centering \large Interaction Transcript for game `codenames', experiment `frequency\_high', episode 6 with reflexion{-}llama{-}3.1{-}8b{-}instruct.
\vspace{24pt}

{ \footnotesize  \setcounter{utterance}{1}
\setlength{\tabcolsep}{0pt}
\begin{supertabular}{c@{$\;$}|p{.15\linewidth}@{}p{.15\linewidth}p{.15\linewidth}p{.15\linewidth}p{.15\linewidth}p{.15\linewidth}}
    \# & $\;$A & \multicolumn{4}{c}{Game Master} & $\;\:$B\\
    \hline

    \theutterance \stepcounter{utterance}  
    & & & \multicolumn{2}{p{0.3\linewidth}}{
        \cellcolor[rgb]{0.9,0.9,0.9}{
            \makecell[{{p{\linewidth}}}]{
                \texttt{\tiny{[GM$|$GM]}}
                {'hidden': {'team': ['road', 'attention', 'bar'], 'innocent': ['school', 'doctor', 'world', 'lord', 'bed', 'stuff'], 'opponent': ['girl', 'evidence'], 'assassin': []}, 'revealed': {'team': {'team': [], 'innocent': [], 'opponent': [], 'assassin': []}, 'opponent': {'team': [], 'innocent': [], 'opponent': [], 'assassin': []}}}
            }
        }
    }
    & & \\ \\

    \theutterance \stepcounter{utterance}  
    & & \multicolumn{4}{p{0.6\linewidth}}{
        \cellcolor[rgb]{0.9,0.9,0.9}{
            \makecell[{{p{\linewidth}}}]{
                \texttt{\tiny{[P1$\langle$GM]}}
                \texttt{You have attempted this same task before and did not succeed. These are the plans you wrote after those attempts.} \\
\texttt{Trial \#1: New plan: } \\
\\ 
\texttt{In my last attempt, I provided clues that were too narrow and specific, which led to my teammate guessing an opponent's word (floor) and another opponent's word (couple) after my clues. This resulted in my team losing the game. Next time, I will provide more general and broader clues that relate to multiple target words, increasing the chances of my teammate guessing all of our team's words before the opposing team does. I will also make sure to avoid providing clues that are too closely related to specific words, and instead focus on more abstract connections between the words. For example, instead of providing a clue like "CLUE: family" that relates to a specific word like "mama", I could provide a clue like "CLUE: family" that relates to multiple words like "mama", "job", and "situation" because they all relate to a family in some way.} \\
\\ 
\texttt{Trial \#2: New plan: } \\
\\ 
\texttt{In my last attempt, I provided a clue that was too specific and only related to two of my team's words ("road" and "bar"), which led to the Game Master rejecting my response. Next time, I will provide a clue that relates to multiple target words, increasing the chances of my teammate guessing all of our team's words before the opposing team does. I will also make sure to provide a clue that is more general and abstract, rather than being closely related to specific words. For example, instead of providing a clue like "CLUE: path" that only relates to "road" and "bar", I could provide a clue like "CLUE: route" that relates to multiple words like "road", "bar", and "attention" because they all relate to a path or a way of getting somewhere.} \\
\\ 
\texttt{Trial \#3: In my last attempt, I provided a clue that was too specific and only related to two of my team's words ("road" and "bar"), which led to the Game Master rejecting my response. Next time, I will provide a clue that relates to all of my team's words, increasing the chances of my teammate guessing all of our team's words before the opposing team does. I will also make sure to provide a clue that is more general and abstract, rather than being closely related to specific words. For example, instead of providing a clue like "CLUE: path" that only relates to "road" and "bar", I could provide a clue like "CLUE: journey" that relates to multiple words like "road", "bar", and "attention" because they all relate to a path or a way of getting somewhere.} \\
\\ 
\\ 
\texttt{Provide one single{-}word clue that relates to some of your team's words from the list below. You can choose to how many and to which words your clue relates to.} \\
\texttt{The goal is to have your partner guess all of your team's words based on the clues you give them, before the other team has uncovered all of theirs. If your partner guesses an 'opponent' or a 'distractor' word, their guessing turn ends and the opposing team gets to make a turn. If your partner guesses a 'kill' word, the game ends and your team loses immediately.} \\
\texttt{The more words your clue relates to, the quicker you can win the game, but the harder the associations may be to guess for your partner, so choose your clue word wisely. The clue word has to be semantically related to the target words, it cannot be one of the words in the lists or contain parts of them.} \\
\\ 
\texttt{Always give your single{-}word clue and your comma{-}separated list of related target words in the following format and make your answers as short as possible, never include any other text than is required in this form:} \\
\\ 
\texttt{CLUE: <WORD>} \\
\texttt{TARGETS: <TARGETS>} \\
\\ 
\texttt{Your team words are: road, attention, bar.} \\
\texttt{Your opponent's team words are: girl, evidence.} \\
\texttt{Distractor words are: school, doctor, world, lord, bed, stuff.} \\
\texttt{The kill words are: .} \\
            }
        }
    }
    & \\ \\

    \theutterance \stepcounter{utterance}  
    & \multicolumn{4}{p{0.6\linewidth}}{
        \cellcolor[rgb]{0.9,0.9,0.9}{
            \makecell[{{p{\linewidth}}}]{
                \texttt{\tiny{[P1$\rangle$GM]}}
                \texttt{CLUE: route} \\
\texttt{TARGETS: road, attention, bar} \\
            }
        }
    }
    & & \\ \\

    \theutterance \stepcounter{utterance}  
    & & & \multicolumn{2}{p{0.3\linewidth}}{
        \cellcolor[rgb]{0.9,0.9,0.9}{
            \makecell[{{p{\linewidth}}}]{
                \texttt{\tiny{[GM$|$GM]}}
                \texttt{route} \\
            }
        }
    }
    & & \\ \\

    \theutterance \stepcounter{utterance}  
    & & & \multicolumn{2}{p{0.3\linewidth}}{
        \cellcolor[rgb]{0.9,0.9,0.9}{
            \makecell[{{p{\linewidth}}}]{
                \texttt{\tiny{[GM$|$GM]}}
                ['road', 'attention', 'bar']
            }
        }
    }
    & & \\ \\

    \theutterance \stepcounter{utterance}  
    & & & \multicolumn{2}{p{0.3\linewidth}}{
        \cellcolor[rgb]{0.9,0.9,0.9}{
            \makecell[{{p{\linewidth}}}]{
                \texttt{\tiny{[GM$|$GM]}}
                {'word': 'road', 'assignment': 'team'}
            }
        }
    }
    & & \\ \\

    \theutterance \stepcounter{utterance}  
    & & & \multicolumn{2}{p{0.3\linewidth}}{
        \cellcolor[rgb]{0.9,0.9,0.9}{
            \makecell[{{p{\linewidth}}}]{
                \texttt{\tiny{[GM$|$GM]}}
                {'word': 'attention', 'assignment': 'team'}
            }
        }
    }
    & & \\ \\

    \theutterance \stepcounter{utterance}  
    & & & \multicolumn{2}{p{0.3\linewidth}}{
        \cellcolor[rgb]{0.9,0.9,0.9}{
            \makecell[{{p{\linewidth}}}]{
                \texttt{\tiny{[GM$|$GM]}}
                {'word': 'bar', 'assignment': 'team'}
            }
        }
    }
    & & \\ \\

    \theutterance \stepcounter{utterance}  
    & & \multicolumn{4}{p{0.6\linewidth}}{
        \cellcolor[rgb]{0.9,0.9,0.9}{
            \makecell[{{p{\linewidth}}}]{
                \texttt{\tiny{[GM$\rangle$P2]}}
                \texttt{You have attempted this same task before and did not succeed. These are the plans you wrote after those attempts.} \\
\texttt{Trial \#1: New plan: I failed to guess the target word because I didn't consider the opponent words as potential targets. In the first round, I correctly guessed the word "horse" but didn't realize it was the target. In the second round, I guessed the word "couple" which was also an opponent word. In the third round, I guessed the word "key" which is closely associated with the word "lock" but not the target. Next time, I will pay closer attention to the opponent words and consider them as potential targets. I will also try to find a word that is closely associated with the target word but not an opponent word. For example, if the target word is "lock", I will try to guess a word that is closely associated with locks, such as "bolt" or "chain".} \\
\\ 
\texttt{Trial \#2: New plan: I failed to provide a well{-}formed and legal response because I included the target word "path" in my guess. In the previous attempts, I correctly guessed words that were closely associated with the target words, but I failed to notice that the target word itself was an option. Next time, I will make sure to include the target word in my guess, in addition to other closely associated words. For example, if the target word is "path", I will guess "GUESS: road, path, trail".} \\
\\ 
\texttt{Trial \#3: New plan: I failed to provide a well{-}formed and legal response because I included the target word "route" in my guess, but I also included another word "road" which is closely associated with the target word. In the previous attempts, I correctly guessed words that were closely associated with the target words, but I failed to notice that the target word itself was an option. Next time, I will make sure to include the target word in my guess, in addition to other closely associated words, but without duplicating the target word. For example, if the target word is "route", I will guess "GUESS: road, path, trail".} \\
\\ 
\\ 
\texttt{Provide a comma{-}separated list of up to 3 words from the following list that best relate or are most closely associated with the word 'route'. Always start your list of guess(es) with 'GUESS: ' and do not include any other text in your answer.} \\
\\ 
\texttt{stuff, doctor, road, girl, bar, school, attention, world, bed, evidence, lord} \\
            }
        }
    }
    & \\ \\

    \theutterance \stepcounter{utterance}  
    & & & \multicolumn{4}{p{0.6\linewidth}}{
        \cellcolor[rgb]{0.9,0.9,0.9}{
            \makecell[{{p{\linewidth}}}]{
                \texttt{\tiny{[GM$\langle$P2]}}
                \texttt{GUESS: road, path, trail} \\
            }
        }
    }
     \\ \\

    \theutterance \stepcounter{utterance}  
    & & & \multicolumn{2}{p{0.3\linewidth}}{
        \cellcolor[rgb]{0.9,0.9,0.9}{
            \makecell[{{p{\linewidth}}}]{
                \texttt{\tiny{[GM$|$GM]}}
                {'player': 'Guesser', 'type': 'guess word is hallucination', 'utterance': 'road, path, trail', 'message': "Guessed word 'path' was not listed, you can only guess words provided in the lists.", 'guess': 'path', 'previous_guesses': [], 'remaining_words': ['stuff', 'doctor', 'road', 'girl', 'bar', 'school', 'attention', 'world', 'bed', 'evidence', 'lord']}
            }
        }
    }
    & & \\ \\

    \theutterance \stepcounter{utterance}  
    & \multicolumn{4}{p{0.6\linewidth}}{
        \cellcolor[rgb]{0.95,0.95,0.95}{
            \makecell[{{p{\linewidth}}}]{
                \texttt{\tiny{[P2$\rightarrow$P1]}}
                \texttt{Post{-}episode Reflexion (not sent to GM):} \\
\texttt{New plan: } \\
\\ 
\texttt{In my last attempt, I provided a clue that was too specific and only related to two of my team's words ("road" and "bar"), which led to the Game Master rejecting my response. Next time, I will provide a clue that relates to all of my team's words, increasing the chances of my teammate guessing all of our team's words before the opposing team does. I will also make sure to provide a clue that is more general and abstract, rather than being closely related to specific words. For example, instead of providing a clue like "CLUE: path" that only relates to "road" and "bar", I could provide a clue like "CLUE: journey" that relates to multiple words like "road", "bar", and "attention" because they all relate to a path or a way of getting somewhere.} \\
            }
        }
    }
    & & \\ \\

    \theutterance \stepcounter{utterance}  
    & & & \multicolumn{4}{p{0.6\linewidth}}{
        \cellcolor[rgb]{0.95,0.95,0.95}{
            \makecell[{{p{\linewidth}}}]{
                \texttt{\tiny{[P1$\rightarrow$P2]}}
                \texttt{Post{-}episode Reflexion (not sent to GM):} \\
\texttt{New plan: I failed to provide a well{-}formed and legal response because I included the target word "route" in my guess, but I also included another word "road" which is closely associated with the target word. In the previous attempts, I correctly guessed words that were closely associated with the target words, but I failed to notice that the target word itself was an option. Next time, I will make sure to include the target word in my guess, in addition to other closely associated words, but without duplicating the target word. However, I also need to pay attention to the format of the response, which requires the guess to be preceded by "GUESS: " and to be a comma{-}separated list of up to 3 words. For example, if the target word is "route", I will guess "GUESS: path, trail, road".} \\
            }
        }
    }
     \\ \\

\end{supertabular}
}

\end{document}
